\section{Introduction}\label{sec:intro}\label{sec:motivation}

The thesis studies how pilots move in cross-country flight: how long they climb in thermals, how
far they glide, how their mean squared displacement grows. These statistics are usually estimated
as if each flight were an independent sample. The assumption fails when pilots fly together: they
share thermals and routes, so their tracks are correlated, and the group may change the behaviour
itself. In soaring birds, social information on the position of thermals changes the gliding
strategy~\cite{williams2018}, and in a migrating flock a few individuals find the thermals while
the others follow them~\cite{flack2018,nagy2018}. To compare group flight with solo flight we must
know, at every instant of every flight, whether the pilot flies in a group and with whom. This
note gives a method that answers, and applies it to the FFVL archive.

\subsection{What flying together means}

The method makes operational an intuitive notion of flying together
(Figure~\ref{fig:requirements}):
\begin{enumerate}[label=(R\arabic*), leftmargin=3.2em]
  \item the pilots take off from the same area, possibly from different takeoffs;
  \item their departures may be staggered, as long as they keep each other in sight;
  \item they follow roughly the same route: nearby places and similar motion; the tracks need
        not pass through the same points;
  \item the composition of the group changes in time: some pilots join, others change route.
\end{enumerate}

\begin{figure}
\centering
\begin{tikzpicture}[x=0.97cm,y=1cm,
    tag/.style={draw=ink2, rounded corners=2pt, fill=white, font=\footnotesize\bfseries, inner sep=2pt}]
  \draw[light, line width=18pt, line cap=round] plot[smooth] coordinates {(2.6,2.35) (6,2.75) (9,2.4) (12,2.4) (14.6,2.4)};
  \draw[ink2, densely dashed] (1.15,2.25) ellipse (0.95 and 1.25);
  \fill[ink] (0.85,2.95) -- ++(-0.13,-0.2) -- ++(0.26,0) -- cycle;
  \fill[ink] (1.4,1.55) -- ++(-0.13,-0.2) -- ++(0.26,0) -- cycle;
  \node[lbl, left] at (0.7,3.0) {A, 11:02};
  \node[lbl, left] at (1.25,1.45) {B, 11:06};
  \draw[track i] plot[smooth] coordinates {(0.85,2.95) (2.6,2.65) (6,3.0) (9,2.65) (10.6,3.0) (12.2,3.9) (13.4,4.3)};
  \draw[track j] plot[smooth] coordinates {(1.4,1.55) (3.2,2.1) (6,2.5) (9,2.15) (12,2.35) (14.6,2.5)};
  \draw[track k] plot[smooth] coordinates {(7.2,0.0) (8.2,1.2) (9.6,1.95) (12,2.15) (14.6,2.25)};
  \node[lbl, blueI, above] at (13.4,4.3) {$i$};
  \node[lbl, orangeJ, right] at (14.6,2.5) {$j$};
  \node[lbl, aqua, right] at (14.6,2.2) {$k$};
  \node[tag] (r1) at (-0.6,0.6) {R1};
  \draw[ink2] (r1) -- (0.5,1.3);
  \node[note, align=left, anchor=north west] at (-1.05,0.3) {same area,\\two takeoffs};
  \node[tag] (r2) at (3.0,0.4) {R2};
  \draw[ink2] (r2) -- (2.0,1.4);
  \node[note, align=left, anchor=north west] at (2.6,0.1) {staggered departures,\\within sight};
  \node[tag] (r3) at (5.6,4.2) {R3};
  \draw[ink2] (r3) -- (6.0,3.15);
  \node[note, align=left, anchor=west] at (6.0,4.2) {nearby places, similar motion};
  \node[tag] (r4) at (10.3,0.5) {R4};
  \draw[ink2] (r4) -- (8.4,1.35);
  \draw[ink2] (r4) -- (11.5,3.45);
  \node[note, align=left, anchor=north west] at (9.9,0.2) {$k$ joins, $i$ changes route};
\end{tikzpicture}
\caption{The four requirements on a schematic map. Pilots $i$ and $j$ take off from two places of
the same area four minutes apart and follow the same corridor (grey) without flying over the same
points; $k$ joins later and $i$ leaves the corridor.}
\label{fig:requirements}
\end{figure}

\subsection{Three difficulties, three ideas}\label{sec:ideas}

Three simple rules miss these requirements, each on one difficulty
(Figure~\ref{fig:approaches}). Each difficulty gives one idea of the method.

\paragraph{1. Together, minutes apart.} Pilots on one route pass the same places minutes apart,
and two minutes are about a kilometre. Rules of strict synchrony (flock~\cite{gudmundsson2006},
convoy~\cite{jeung2008}, swarm~\cite{li2010}) ask for proximity at the same instant and never see
them together (Figure~\ref{fig:approaches}b); the swarm admits non-consecutive instants, but no
delay between members. \emph{Idea:} two pilots are in contact when one of them is where the other
was, or will be, at most a few minutes earlier or later. This is the following with a local delay
of Li, Wu and Crofoot~\cite{li2013}, which we adapt to soaring in three ways: the tolerance is a
flat box, wide horizontally and narrow in height, like a group; headings are compared only in
glide, where they show the route; the takeoff area, where everyone is close to everyone, is
left out (Section~\ref{sec:contact}).

\paragraph{2. Not transitive, and changing in time.} If 1 flies with 8 early in the afternoon and
8 with 14 later, 1 and 14 may never have met; and pilots join and leave all the time (R4). A
window on the departures, such as ``departures within $T$ minutes'', makes groups that depend on
the reference pilot (Figure~\ref{fig:approaches}c). \emph{Idea:} groups are defined instant by
instant, as the connected components of the graph of contacts: each component, with its edges,
is a \emph{configuration}. A fixed set of pilots that stays one component is a \emph{group},
and a chain of groups that share members is an \emph{episode}, after the grouping structure of Buchin et al.~\cite{buchin2015}. Episodes tie the
flights into \emph{dependency blocks}, which the statistics resample together
(Section~\ref{sec:groups}).

\paragraph{3. Chance meetings.} Routes follow the terrain and many thermals are anchored to it:
independent pilots from the same site fly the same ridges and climb in the same thermals.
Geometric similarity of the tracks (Hausdorff, Fréchet, unconstrained DTW; survey
in~\cite{su2020}) finds such tracks similar even hours apart (Figure~\ref{fig:approaches}a);
unconstrained DTW also warps time without limit and loses the delay that (R2) bounds.
\emph{Idea:} chance is measured on the data. Copies of the flights shifted in time keep the routes
and lose the synchrony, so their contacts are chance. A contact counts only if it lasts at least
$\tau^*$, chosen so that at most one contact in ten is chance: $\tau^*=\tauStar$~min on the
archive (Section~\ref{sec:calibration}).

\begin{figure}
\centering
\begin{tikzpicture}[x=1cm,y=1cm]
  % (a) geometric similarity
  \begin{scope}[shift={(0,0)}]
    \node[panel] at (0,3.3) {(a) Geometric similarity};
    \draw[muted, line width=0.5pt] plot[smooth] coordinates {(0,2.6) (1.5,2.95) (3,2.7) (4.6,3.0)};
    \draw[muted, line width=0.5pt] plot[smooth] coordinates {(0,0.6) (1.5,0.95) (3,0.7) (4.6,1.0)};
    \node[note, anchor=north] at (3.9,2.75) {valley sides};
    \draw[track i] plot[smooth] coordinates {(0.1,1.6) (1.2,2.0) (2.4,1.7) (3.5,2.05) (4.5,1.85)};
    \draw[track j] plot[smooth] coordinates {(0.1,1.45) (1.2,1.83) (2.4,1.55) (3.5,1.88) (4.5,1.7)};
    \node[lbl, blueI, above] at (0.6,1.85) {$i$, 10:00};
    \node[lbl, orangeJ, below] at (0.75,1.4) {$j$, 14:00};
    \node[note, align=left, anchor=north west] at (0,0.35) {same shape,\\four hours apart};
  \end{scope}
  % (b) strict synchrony
  \begin{scope}[shift={(5.3,0)}]
    \node[panel] at (0,3.3) {(b) Strict synchrony};
    \draw[light, line width=6pt] (0.1,1.5) -- (4.6,1.5);
    \draw[track i] (0.1,1.5) -- (3.4,1.5);
    \draw[track j] (0.1,1.5) -- (1.4,1.5);
    \node[pt i] at (3.4,1.5) {};
    \node[pt j] at (1.4,1.5) {};
    \draw[tol] (3.4,1.5) circle (0.7);
    \draw[<->, ink2] (1.4,2.35) -- node[above, lbl] {2 min $\approx$ 1 km} (3.4,2.35);
    \node[lbl, orangeJ, below=3pt] at (1.4,1.5) {$j$};
    \node[lbl, blueI, below=3pt] at (3.4,1.5) {$i$};
    \node[note, align=left, anchor=north west] at (0,0.35) {together, never close\\at the same instant};
  \end{scope}
  % (c) window on departures
  \begin{scope}[shift={(10.6,0)}]
    \node[panel] at (0,3.3) {(c) Window on departures};
    \draw[->, ink2] (0,1.0) -- (4.9,1.0) node[right, lbl, ink] {$t$};
    \foreach \x/\f/\h in {0.4/1/11:00, 2.0/8/11:08, 3.6/14/11:16} {
      \fill[ink] (\x,1.0) circle (1.6pt);
      \node[lbl, below] at (\x,0.95) {\h};
      \node[lbl, above] at (\x,1.05) {\f};
    }
    \draw[ink2] (0.4,1.6) -- (2.0,1.6) node[midway, above, lbl] {\checkmark};
    \draw[ink2] (2.0,1.95) -- (3.6,1.95) node[midway, above, lbl] {\checkmark};
    \draw[ink2] (0.4,2.45) -- (3.6,2.45) node[midway, above, lbl] {$\times$};
    \node[note, align=left, anchor=north west] at (0,0.35) {window $T = 10$ min: the group\\depends on the reference pilot};
  \end{scope}
\end{tikzpicture}
\caption{Why three immediate rules fail. (a) Two independent pilots on the same valley produce
similar tracks hours apart. (b) Two pilots on the same route two minutes apart never enter each
other's proximity disc. (c) With a window of 10 minutes, 1--8 and 8--14 pass and 1--14 fails.}
\label{fig:approaches}
\end{figure}

\subsection{The method in six steps}

The three ideas chain into six steps, from two tracks to the units of the statistics
(Figure~\ref{fig:overview}); Table~\ref{tab:glossary} names each object in plain words, with the
section that defines it. Figures~\ref{fig:contact} and~\ref{fig:day}, marked
\animmark, are animated: in Adobe Acrobat Reader, or another viewer that runs the JavaScript of
PDF files, the buttons below them play the animation; other viewers show the figure as printed.
The same animations are GIF files in the folder \texttt{animations/}.

\begin{figure}
\centering
\begin{tikzpicture}[x=1cm,y=1cm]
  % 1: two tracks
  \begin{scope}[shift={(0,0)}]
    \node[panel] at (0,3.25) {1\enspace Two tracks};
    \draw[track i] plot[smooth] coordinates {(0.2,1.3) (1.2,1.9) (2.4,1.8) (3.4,2.45) (4.4,2.8)};
    \draw[track j] plot[smooth] coordinates {(0.2,0.85) (1.2,1.45) (2.4,1.35) (3.4,2.0) (4.4,2.35)};
    \node[pt i] (xi) at (3.4,2.45) {};
    \node[hollow j] (xj) at (1.2,1.45) {};
    \node[lbl, blueI, above=1pt] at (xi) {$\x_i(t)$};
    \node[lbl, orangeJ, below=2pt] at (xj) {$\x_j(t)$};
    \node[note, align=left, anchor=north west] at (0.1,0.55) {$j$ flies a similar route,\\a little later};
  \end{scope}
  % 2: correspondences
  \begin{scope}[shift={(5.45,0)}]
    \node[panel] at (0,3.25) {2\enspace Correspondences $\M_{ij}$};
    \fill[light] (0,0) -- (0.6,0) -- (2.6,2.0) -- (2.6,2.6) -- (2.0,2.6) -- (0,0.6) -- cycle;
    \draw[muted, densely dashed] (0,0) -- (2.6,2.6);
    \fill[aqua, rotate around={45:(1.15,1.5)}] (1.15,1.5) ellipse (0.78 and 0.12);
    \draw[->, ink2] (0,0) -- (2.85,0) node[right, lbl, ink] {$t$};
    \draw[->, ink2] (0,0) -- (0,2.85) node[above, lbl, ink] {$t'$};
    \node[note, align=left, anchor=west] at (2.75,1.55) {pairs of\\instants;\\above the\\diagonal:\\$j$ is later};
  \end{scope}
  % 3: contact
  \begin{scope}[shift={(10.9,0)}]
    \node[panel] at (0,3.25) {3\enspace Contact $a_{ij}(t)$};
    \draw[->, ink2] (0,1.0) -- (4.5,1.0) node[right, lbl, ink] {$t$};
    \fill[light] (0,1.3) rectangle (4.3,1.8);
    \foreach \a/\b in {0.3/1.15, 1.3/2.2, 2.75/2.9, 3.55/3.66} \fill[aqua] (\a,1.3) rectangle (\b,1.8);
    \node[note, align=left, anchor=north west] at (0,0.75) {$1$ when one of the two pilots\\is on the path of the other};
  \end{scope}
  % 4: pair intervals
  \begin{scope}[shift={(0,-4.2)}]
    \node[panel] at (0,3.25) {4\enspace Pair intervals};
    \fill[light] (0,2.1) rectangle (4.3,2.5);
    \foreach \a/\b in {0.3/1.15, 1.3/2.2, 2.75/2.9, 3.55/3.66} \fill[aqua!45] (\a,2.1) rectangle (\b,2.5);
    \fill[light] (0,1.2) rectangle (4.3,1.6);
    \fill[aqua] (0.3,1.2) rectangle (2.2,1.6);
    \draw[->, ink2] (0,0.9) -- (4.5,0.9) node[right, lbl, ink] {$t$};
    \node[lbl, below] at (0.3,0.9) {$t_{\rm on}$};
    \node[lbl, below] at (2.2,0.9) {$t_{\rm off}$};
    \node[note, align=left, anchor=north west] at (0,0.35) {short gaps filled, short pieces dropped};
  \end{scope}
  % 5: groups at each instant
  \begin{scope}[shift={(5.45,-4.2)}]
    \node[panel] at (0,3.25) {5\enspace Configurations, groups};
    \coordinate (n1) at (0.45,1.9); \coordinate (n8) at (1.4,2.55); \coordinate (n22) at (1.4,1.25);
    \coordinate (n14) at (2.45,1.9); \coordinate (n5) at (0.45,0.75);
    \draw[edge] (n1) -- (n8) -- (n22);
    \node[inE1] at (n1) {1}; \node[inE1] at (n8) {8}; \node[inE1] at (n22) {22};
    \node[flight] at (n14) {14}; \node[flight] at (n5) {5};
    \node[note, align=left, anchor=west] at (2.7,1.75) {linked pilots:\\configuration;\\same members\\in time: group};
  \end{scope}
  % 6: episodes and blocks
  \begin{scope}[shift={(10.9,-4.2)}]
    \node[panel] at (0,3.25) {6\enspace Episodes and blocks};
    \foreach \a/\b/\n in {0.1/1.15/1, 1.2/1.75/2, 1.8/2.6/3} {
      \fill[blueI] (\a,1.9) rectangle (\b,2.4);
      \node[font=\footnotesize, text=white] at ({(\a+\b)/2},2.15) {$g_\n$};
    }
    \fill[orangeJ] (3.3,1.9) rectangle (4.1,2.4);
    \node[font=\footnotesize, text=white] at (3.7,2.15) {$g_4$};
    \draw[decorate, decoration={brace, mirror, amplitude=4pt}, ink2] (0.1,1.8) -- (2.6,1.8)
      node[midway, below=4pt, lbl, ink] {episode E1};
    \draw[decorate, decoration={brace, mirror, amplitude=4pt}, ink2] (3.3,1.8) -- (4.1,1.8)
      node[midway, below=4pt, lbl, ink] {E2};
    \node[note, align=left, anchor=north west] at (0,0.75) {block: the flights tied by episodes,\\resampled together};
  \end{scope}
  % reading order
  \foreach \y in {1.6, -2.6} {
    \draw[->, muted, line width=0.8pt] (4.75,\y) -- (5.3,\y);
    \draw[->, muted, line width=0.8pt] (10.2,\y) -- (10.75,\y);
  }
\end{tikzpicture}
\caption{The method in six steps. Panels 1--3: the pair contact (Section~\ref{sec:contact}).
Panels 4--6: from pairs to groups (Section~\ref{sec:groups}). Flights 1, 8, 22, 14 and 5 are the
example day of Figure~\ref{fig:day}.}
\label{fig:overview}
\end{figure}

\begin{table}
\centering
\small
\begin{tabularx}{\textwidth}{@{}l X l l@{}}
\toprule
Object & In words & Symbol & Section \\
\midrule
Departure & the first fix farther than $R_\dep$ from the takeoff & $k_i^\dep$ & \ref{sec:local} \\
Correspondence & two instants of the same clock, $t$ for $i$ and $t'$ for $j$, at most $\delta$ apart, at which the two pilots pass the same place within the tolerances & $(t,t')\in\M_{ij}$ & \ref{sec:correspondences} \\
Contact & at instant $t$, one of the two pilots is on the path of the other & $a_{ij}(t)=1$ & \ref{sec:correspondences} \\
Pair interval & a stretch of contact, with short gaps filled and short pieces dropped & $[t_{\rm on},t_{\rm off})$ & \ref{sec:pair-intervals} \\
Configuration & the pilots joined to each other by chains of contacts at instant $t$, with their edges; a new edge gives a new configuration & component of $G(t)$ & \ref{sec:configurations} \\
Group & a fixed set of flights that stays one component over a maximal interval & $g$ & \ref{sec:configurations} \\
Membership & the label of flight $i$ at instant $t$: local phase, solo, or its group & $\gamma_i(t)$ & \ref{sec:configurations} \\
Episode & a chain of groups, each starting when the previous one ends and sharing a member with it & & \ref{sec:episodes} \\
Dependency block & the flights tied together by episodes; the bootstrap resamples them together & & \ref{sec:blocks} \\
\bottomrule
\end{tabularx}
\caption{A minimal dictionary of the objects of the method.}
\label{tab:glossary}
\end{table}

\subsection{What we take and what is new}

Each ingredient of the method has a precedent (Appendix~\ref{sec:literature},
Table~\ref{tab:origin}). The pair contact is the following with a local delay of Li et
al.~\cite{li2013}; the group structure is that of Buchin et al.~\cite{buchin2015}; groups built
from delayed relations exist in mFLICA~\cite{amornbunchornvej2021}, directed, with leaders and on
sliding windows; the time shift of the null model comes from Spiegel et al.~\cite{spiegel2016}.
What we did not find in the literature is:
\begin{enumerate}[label=(\roman*)]
  \item the Reeb structure of~\cite{buchin2015} built on contacts with a local delay instead of
        simultaneous proximity;
  \item the adaptations to thermal flight: a vertical tolerance that separates pilots circling in
        the same thermal at distant heights, and headings compared only in glide
        (Section~\ref{sec:roles});
  \item the application to a free-flight dataset of this size, with thresholds calibrated on a
        null model.
\end{enumerate}
The contribution is therefore a physically motivated adaptation and a scalable implementation,
more than a new method in itself, and the thesis should present it in these terms.

\subsection{Plan of the note}

Section~\ref{sec:data} fixes the data and the notation. Section~\ref{sec:contact} defines the
contact of two pilots, Section~\ref{sec:groups} builds the groups from it, and
Section~\ref{sec:calibration} defines chance and the minimum duration $\tau^*$.
Section~\ref{sec:tuning} sets the parameters: measured on the data where the data fix them,
declared as definitions where they do not. Section~\ref{sec:use-limits} says how the labels serve the comparison of group and solo
flight, and lists the choices to declare and the limits. What the groups are, what holds them
together, who leads them and with whom pilots fly is in the note \emph{The structure of groups in
free flight}. The appendices hold the literature
(Appendix~\ref{sec:literature}), the delay of a contact that says who follows whom
(Appendix~\ref{sec:lag}), the algorithm, with the proof that its filter loses no contact
(Appendix~\ref{sec:algorithm}), the output tables (Appendix~\ref{sec:outputs}), why the two halves
of the contact differ (Appendix~\ref{app:asymmetry}), the terms of graph theory
(Appendix~\ref{app:graphs}) and further evidence (Appendix~\ref{app:evidence}).
